\chapter{技术栈与项目结构}

\section{为什么选择 Bun}

Claude Code 选择 \href{https://bun.sh}{Bun} 作为运行时而非 Node.js，这是一个值得深思的技术决策。

\subsection{Bun 的关键优势}

\begin{longtable}{llll}
\toprule
特性 & Node.js & Bun & Claude Code 的受益点 \\
\midrule
启动速度 & \textasciitilde{}40ms & \textasciitilde{}6ms & CLI 工具对冷启动极度敏感 \\
TypeScript 支持 & 需要 ts-node/tsx & 原生支持 & 无需编译步骤即可运行 .ts/.tsx \\
内置 bundler & 无 & 有 & 支持 \code{bun:bundle} Feature Flag 系统 \\
包管理器 & npm/yarn/pnpm & 内置 (bun install) & 更快的依赖安装 \\
文件系统操作 & 标准库 & 优化的 I/O & 大量文件读写操作受益 \\
\bottomrule
\end{longtable}

\subsection{\code{bun:bundle} 虚拟模块}

Claude Code 利用 Bun 独有的 \code{bun:bundle} 虚拟模块实现编译时功能开关：

\begin{lstlisting}[language=TypeScript,breaklines=true]
// 生产环境：Bun bundler 在编译时静态替换 feature() 调用
import { feature } from 'bun:bundle'

// 如果 VOICE_MODE 未启用，以下代码会在编译时被完全移除（死代码消除）
const voiceCommand = feature('VOICE_MODE')
  ? require('./commands/voice/index.js').default
  : null
\end{lstlisting}

在开发环境中，通过 \code{plugins/bunBundleDev.ts} 插件提供运行时 polyfill：

\begin{lstlisting}[language=TypeScript,breaklines=true]
// plugins/bunBundleDev.ts
import { plugin } from 'bun'

const enabledFlags = new Set(
  (process.env.FEATURE_FLAGS ?? '').split(',').filter(Boolean),
)

plugin({
  name: 'bun-bundle-dev',
  setup(build) {
    build.module('bun:bundle', () => ({
      exports: {
        feature(flag: string): boolean {
          return enabledFlags.has(flag)
        },
      },
      loader: 'object',
    }))
  },
})
\end{lstlisting}

开发时通过环境变量启用特定功能：

\begin{lstlisting}[language=bash,breaklines=true]
FEATURE_FLAGS=KAIROS,VOICE_MODE bun run src/main.tsx
\end{lstlisting}

\section{TypeScript 配置}

Claude Code 使用严格模式的 TypeScript：

\begin{lstlisting}[breaklines=true]
{
  "compilerOptions": {
    "target": "ESNext",
    "module": "ESNext",
    "moduleResolution": "bundler",
    "jsx": "react-jsx",
    "strict": true,
    "esModuleInterop": true,
    "skipLibCheck": true,
    "forceConsistentCasingInFileNames": true,
    "resolveJsonModule": true,
    "baseUrl": ".",
    "paths": {
      "src/*": ["./src/*"],
      "bun:bundle": ["./src/types/bun-bundle.d.ts"]
    },
    "types": ["bun-types"],
    "lib": ["ESNext"],
    "allowImportingTsExtensions": true,
    "noEmit": true
  }
}
\end{lstlisting}

关键配置点：

\begin{enumerate}[nosep]
  \item \textbf{\code{moduleResolution: "bundler"}}：使用 Bun 的模块解析策略。
  \item \textbf{\code{jsx: "react-jsx"}}：支持 JSX 语法，用于 Ink 终端 UI 组件。
  \item \textbf{路径别名}：\code{src/*} 映射允许使用绝对导入路径（如 \code{src/bootstrap/state.js}）。
  \item \textbf{\code{bun:bundle} 类型映射}：将虚拟模块映射到本地类型定义文件。
\end{enumerate}

\section{构建时宏定义}

通过 \code{bunfig.toml} 定义构建时注入的全局常量：

\begin{lstlisting}[breaklines=true]
[run]
preload = ["./plugins/bunBundleDev.ts"]

[define]
"MACRO.VERSION" = '"1.0.0-dev"'
"MACRO.BUILD_TIME" = '""'
"MACRO.PACKAGE_URL" = '"@anthropic-ai/claude-code"'
"MACRO.NATIVE_PACKAGE_URL" = 'undefined'
"MACRO.FEEDBACK_CHANNEL" = '"https://github.com/anthropics/claude-code/issues"'
"MACRO.ISSUES_EXPLAINER" = '"file an issue at ..."'
"MACRO.VERSION_CHANGELOG" = '""'
\end{lstlisting}

这些 \code{MACRO.*} 常量在代码中直接作为全局变量使用，Bun 在构建/运行时自动替换。

\section{依赖分析}

\subsection{核心依赖}

\begin{longtable}{lll}
\toprule
分类 & 包名 & 用途 \\
\midrule
\textbf{AI/API} & \code{@anthropic-ai/sdk} & Anthropic API 客户端 \\
 & \code{@anthropic-ai/claude-agent-sdk} & Agent SDK \\
 & \code{@aws-sdk/client-bedrock-runtime} & AWS Bedrock 支持 \\
 & \code{google-auth-library} & Google Cloud 认证 \\
\textbf{终端 UI} & \code{react} (canary) & React 19 Canary 版本 \\
 & \code{react-reconciler} (canary) & 自定义 React 渲染器 \\
 & \code{chalk} & 终端着色 \\
 & \code{cli-boxes} & 终端边框 \\
 & \code{figures} & Unicode 符号 \\
\textbf{CLI} & \code{commander} & 命令解析 \\
 & \code{@commander-js/extra-typings} & Commander 增强类型 \\
\textbf{协议} & \code{@modelcontextprotocol/sdk} & MCP 协议实现 \\
 & \code{vscode-jsonrpc} & JSON-RPC 通信 \\
 & \code{vscode-languageserver-protocol} & LSP 协议 \\
\textbf{数据处理} & \code{zod} & 运行时模式验证 \\
 & \code{ajv} & JSON Schema 验证 \\
 & \code{diff} & 文本差异比较 \\
 & \code{marked} & Markdown 解析 \\
\textbf{工具} & \code{execa} & 子进程管理 \\
 & \code{chokidar} & 文件监听 \\
 & \code{fuse.js} & 模糊搜索 \\
 & \code{lodash-es} & 实用函数库 \\
\textbf{网络} & \code{undici} & HTTP 客户端 \\
 & \code{ws} & WebSocket \\
 & \code{https-proxy-agent} & 代理支持 \\
\textbf{安全} & \code{xss} & XSS 过滤 \\
 & \code{proper-lockfile} & 文件锁 \\
\textbf{遥测} & \code{@opentelemetry/*} & OpenTelemetry 全家桶 \\
 & \code{@growthbook/growthbook} & 功能开关 \\
\bottomrule
\end{longtable}

\subsection{特殊依赖：React Canary}

值得注意的是，Claude Code 使用的是 React \textbf{19.3.0-canary} 版本，而非稳定版。这是因为 Ink（终端 UI 框架）需要与 \code{react-reconciler} 的 canary 版本匹配。

\subsection{内部 Stub 包}

被替换的 Anthropic 内部包（使用 \code{file:stubs/...} 引用）：

\begin{lstlisting}[breaklines=true]
{
  "@ant/claude-for-chrome-mcp": "file:stubs/@ant/claude-for-chrome-mcp",
  "@ant/computer-use-input": "file:stubs/@ant/computer-use-input",
  "@ant/computer-use-mcp": "file:stubs/@ant/computer-use-mcp",
  "@ant/computer-use-swift": "file:stubs/@ant/computer-use-swift",
  "@anthropic-ai/mcpb": "file:stubs/@anthropic-ai/mcpb",
  "@anthropic-ai/sandbox-runtime": "file:stubs/anthropic-ai/sandbox-runtime"
}
\end{lstlisting}

这些包在内部用于：
\begin{itemize}[nosep]
  \item \textbf{Chrome 浏览器控制} (\code{claude-for-chrome-mcp})
  \item \textbf{计算机使用} (\code{computer-use-*})：屏幕截图、鼠标/键盘操控
  \item \textbf{沙箱运行时} (\code{sandbox-runtime})：安全隔离执行环境
\end{itemize}

\section{目录结构详解}

\begin{lstlisting}[language=TypeScript,breaklines=true]
src/
├── main.tsx                 # 入口文件：Commander.js CLI 解析 + React/Ink 初始化
├── commands.ts              # 命令注册表：管理所有斜杠命令
├── tools.ts                 # 工具注册表：管理所有 AI 可调用工具
├── Tool.ts                  # 工具类型定义：接口、权限模型、进度类型
├── QueryEngine.ts           # 核心查询引擎：LLM API 调用、流式响应、工具循环
├── query.ts                 # 查询循环：消息处理、自动压缩、Token 管理
├── context.ts               # 上下文收集：System/User Context 生成
├── cost-tracker.ts          # 成本追踪：Token 计数、费用计算
│
├── commands/                # 斜杠命令实现 (~50 个命令)
│   ├── commit.ts            #   /commit - Git 提交
│   ├── review.ts            #   /review - 代码审查
│   ├── compact/             #   /compact - 上下文压缩
│   ├── config/              #   /config - 设置管理
│   ├── doctor/              #   /doctor - 环境诊断
│   ├── memory/              #   /memory - 持久记忆管理
│   ├── mcp/                 #   /mcp - MCP 服务器管理
│   ├── skills/              #   /skills - 技能管理
│   ├── tasks/               #   /tasks - 任务管理
│   ├── vim/                 #   /vim - Vim 模式
│   └── ...
│
├── tools/                   # 工具实现 (~40 个工具)
│   ├── AgentTool/           #   子代理生成
│   ├── BashTool/            #   Shell 命令执行
│   ├── FileReadTool/        #   文件读取
│   ├── FileWriteTool/       #   文件创建/覆盖
│   ├── FileEditTool/        #   文件部分编辑
│   ├── GlobTool/            #   文件模式匹配
│   ├── GrepTool/            #   ripgrep 内容搜索
│   ├── WebFetchTool/        #   URL 内容获取
│   ├── WebSearchTool/       #   网页搜索
│   ├── SkillTool/           #   技能执行
│   ├── MCPTool/             #   MCP 工具调用
│   ├── LSPTool/             #   语言服务器协议
│   ├── TeamCreateTool/      #   团队代理创建
│   ├── SendMessageTool/     #   代理间消息
│   ├── EnterPlanModeTool/   #   进入计划模式
│   ├── EnterWorktreeTool/   #   Git Worktree 隔离
│   └── ...
│
├── components/              # Ink UI 组件 (~140 个)
│   ├── REPL.tsx             #   主 REPL 界面
│   ├── MessageSelector.tsx  #   消息选择器
│   ├── AgentProgressLine.tsx #  代理进度显示
│   ├── Spinner.tsx          #   加载动画
│   └── ...
│
├── hooks/                   # React Hooks
│   ├── toolPermission/      #   工具权限检查
│   ├── fileSuggestions.ts   #   文件建议
│   ├── useCanUseTool.ts     #   工具使用权限
│   └── ...
│
├── services/                # 外部服务集成
│   ├── api/                 #   Anthropic API 客户端
│   ├── mcp/                 #   MCP 服务器管理
│   ├── oauth/               #   OAuth 认证
│   ├── lsp/                 #   LSP 管理器
│   ├── compact/             #   上下文压缩服务
│   ├── analytics/           #   GrowthBook 分析
│   ├── plugins/             #   插件加载
│   ├── policyLimits/        #   策略限制
│   ├── extractMemories/     #   自动记忆提取
│   └── ...
│
├── bridge/                  # IDE 集成桥
│   ├── bridgeMain.ts        #   桥主循环
│   ├── bridgeMessaging.ts   #   消息协议
│   └── ...
│
├── coordinator/             # 多代理协调器
│   └── coordinatorMode.ts   #   协调器模式
│
├── state/                   # 状态管理
│   ├── AppStateStore.tsx    #   全局应用状态
│   ├── store.ts             #   状态存储
│   └── onChangeAppState.ts  #   状态变更监听
│
├── plugins/                 # 插件系统
├── skills/                  # 技能系统
├── memdir/                  # 持久记忆目录
├── tasks/                   # 任务管理
├── voice/                   # 语音输入
├── vim/                     # Vim 模式
├── remote/                  # 远程会话
├── server/                  # 服务器模式
├── schemas/                 # Zod 配置模式
├── migrations/              # 配置迁移
├── entrypoints/             # 入口初始化逻辑
├── types/                   # 类型定义
├── utils/                   # 工具函数
├── constants/               # 常量定义
├── ink/                     # Ink 渲染器封装
├── buddy/                   # 伴侣精灵（彩蛋）
└── upstreamproxy/           # 代理配置
\end{lstlisting}

\section{设计启示}

\subsection{选择合适的运行时}

Claude Code 选择 Bun 而非 Node.js 的决策展示了一个重要原则：\textbf{CLI 工具应该为冷启动时间进行优化}。Bun 的 6ms 启动时间 vs Node.js 的 40ms，对于频繁调用的 CLI 工具来说影响显著。

\subsection{虚拟模块实现编译时优化}

\code{bun:bundle} 的 Feature Flag 系统展示了一种优雅的\textbf{死代码消除}方案。未启用的功能在编译时被完全移除，减小了最终产物的体积，同时避免了运行时检查的开销。

\subsection{内部包的 Stub 策略}

对于无法公开的内部依赖，使用 \code{file:stubs/} 替换为空实现是一种实用的解耦策略，使得核心代码可以在缺失内部依赖的情况下独立运行。

\bigskip\noindent\rule{\textwidth}{0.4pt}\bigskip

\textbf{下一章}：\href{./chapter-03-startup.md}{启动流程与初始化}
